\section*{Capítulo \ref{ch:b3:glaciology} --- Glaciologia}

\begin{solution}{exo:b3:glaciology:1}
$\tau_b = \rho gH\sin\alpha = 917\times9.81\times250\times\sin\ang{3} =
\qty{118}{kPa}$, cerca de \qty{1.2}{bar}: a ordem de grandeza das tensões
sob geleiras reais.
\end{solution}

\begin{solution}{exo:b3:glaciology:2}
(a)~$\qty{2.9e15}{m^3}\times917 = \qty{2.66e18}{kg} = \qty{2.66e6}{Gt}$;
dividido por \qty{360}{Gt/mm} dá \qty{7390}{mm}, \qty{7.4}{m}, o valor citado.
(b)~$\qty{3.0e16}{m^3}\times917/360 = \qty{76}{m}$, contra \qty{58}{m}.
O gelo da Groenlândia repousa em sua maior parte sobre rocha acima do nível do mar, por isso todo ele
acrescentaria água ao oceano. Grande parte do gelo da Antártida repousa sobre um leito abaixo do
nível do mar: esse gelo já desloca água do mar, e derretê-lo apenas repõe a água que ele desloca;
as plataformas de gelo flutuantes não acrescentam nada.
\end{solution}

\begin{solution}{exo:b3:glaciology:3}
Ganho $\num{6e6}\times1.2 = \qty{7.2e6}{m^3}$, perda $\num{4e6}\times2.1 =
\qty{8.4e6}{m^3}$ de água: saldo líquido de \qty{-1.2e6}{m^3}, com balanço específico médio
de $-1.2\times10^6/10^7 = \qty{-0.12}{m}$ equivalente em água. A
razão da área de acumulação é $6/10 = 0.6$. A geleira está perdendo massa:
encolhendo.
\end{solution}

\begin{solution}{exo:b3:glaciology:4}
$1 - 820/917 = \num{0.106}$ e $1 - 830/917 = \num{0.095}$: cerca de um décimo
do volume ainda é espaço poroso no momento do \omterm{def:b3:glaciology:firn}{fechamento dos poros}.
\end{solution}

\begin{solution}{exo:b3:glaciology:5}
$\tau_b = 917\times9.81\times300\times\sin\ang{3} = \qty{141.2}{kPa}$. Pelo
teorema da laje, $u_s = \tfrac{2A}{4}\tau_b^3H = 0.5\times\num{1.43e-16}
\times(\num{1.412e5})^3\times300 = \qty{60}{m/yr}$. Com o dobro da espessura: $u_s
\propto H^{n+1} = H^4$, dezesseis vezes mais rápida. Com metade da inclinação: $u_s \propto
(\sin\alpha)^3$, oito vezes mais lenta.
\end{solution}

\begin{solution}{exo:b3:glaciology:6}
A \qty{-25}{\celsius} (\qty{248.15}{K}), $A = \num{3.61e-13}\exp(-60000/
(8.314\times248.15)) = \qty{8.5e-26}{Pa^{-3}.s^{-1}}$, ou
\qty{2.7e-18}{Pa^{-3}.yr^{-1}}. A \qty{-10}{\celsius}, $\num{3.61e-13}
\exp(-27.42) = \qty{4.4e-25}{Pa^{-3}.s^{-1}} = \qty{1.4e-17}{Pa^{-3}.yr^{-1}}$.
A \qty{0}{\celsius} (ramo quente) \qty{1.4e-16}{Pa^{-3}.yr^{-1}}. Sob
as mesmas tensões, as velocidades estão na razão dos \omterm{def:b3:glaciology:glen}{fatores de taxa}:
$\num{1.43e-16}/\num{2.67e-18} = 54$.
\end{solution}

\begin{solution}{exo:b3:glaciology:7}
$h = \sqrt{2\times10^5\times10^6/(917\times9.81)} = \qty{4.7}{km}$, próxima dos
\qty{4.9}{km} da Antártida: o modelo plástico obtém a espessura de um
manto de gelo apenas a partir de sua largura.
\end{solution}

\begin{solution}{exo:b3:glaciology:8}
$\tau = 200/5 = \num{40}\text{ anos}$. O comprimento continuaria a se ajustar por
várias décadas, alguns tempos de resposta, depois que o aquecimento cessou: as
geleiras de hoje ainda não acompanharam o clima atual.
\end{solution}

\begin{solution}{exo:b3:glaciology:9}
(a)~$h = \sqrt{2\times2.03\times20\times\num{8.64e6}/(917\times\num{3.34e5})}
= \qty{1.51}{m}$. Como $h^2 \propto t$, atingir \qty{3}{m} exige $100\times
(3/1.51)^2 = 393$ dias, mais que um inverno. (b)~Gelo muito mais espesso que
\qty{3}{m} não pode se formar por congelamento num único inverno: é gelo que
sobreviveu a vários, ou gelo empilhado pelo vento e pelas correntes em cristas de pressão.
\end{solution}

\begin{solution}{exo:b3:glaciology:10}
(a)~$H = (1025/917)\times800 = \qty{894}{m}$. (b)~Com \qty{900}{m}, a
espessura da \omterm{def:b3:glaciology:grounding-line}{linha de encalhe} é \qty{1006}{m}, $9/8$ da anterior; o fluxo
aumenta em $(9/8)^5 = 1.8$. (c)~Conforme a linha recua por um leito que se aprofunda
para o interior, o gelo na linha é mais espesso, descarrega mais, afina o gelo
atrás, e a linha recua ainda mais: uma retroalimentação positiva. Uma plataforma de gelo
que escora o gelo apoiado no leito reduz o fluxo na \omterm{def:b3:glaciology:grounding-line}{linha de encalhe} e pode
mantê-la no lugar mesmo em uma inclinação descendente.
\end{solution}

\begin{solution}{exo:b3:glaciology:11}
Com $h(r) = c\sqrt{R - r}$, $c = \sqrt{2\tau_0/(\rho g)}$, $V =
\int_0^R2\pi rc\sqrt{R - r}\,\mathrm dr$. Substituindo $s = R - r$,
\[
\int_0^R(R - s)\sqrt s\,\mathrm ds = \tfrac23R^{5/2} - \tfrac25R^{5/2} =
\tfrac4{15}R^{5/2},
\]
logo $V = \frac{8\pi}{15}cR^{5/2}$. Ao dobrar $R$
multiplica $V$ e o \omterm{def:b3:glaciology:sle}{equivalente em nível do mar} por $2^{5/2} = 5.7$.
\end{solution}

\begin{solution}{exo:b3:glaciology:12}
(a)~Março: $14.13/16.34 - 1 = \qty{-13.5}{\percent}$; setembro: $4.81/7.05
- 1 = \qty{-31.8}{\percent}$. (b)~Em março, a borda do gelo é mantida em grande parte pelas
costas que cercam o oceano Ártico, e o recongelamento a cada inverno
volta a preencher a bacia; o gelo mais fino dos últimos anos derrete
por completo a cada verão. E a retroalimentação gelo--albedo atua no verão: a água aberta
absorve cerca de \qty{94}{\percent} da luz solar, o gelo reflete
\qtyrange{50}{70}{\percent}, de modo que cada perda de gelo de verão aquece a água e
derrete mais.
\end{solution}

\begin{solution}{pb:b3:glaciology:1}
\textbf{1.} A coluna acima do leito, de peso $\rho gH$ por unidade de área, é
empurrada encosta abaixo pela componente $\rho gH\sin\alpha$, equilibrada pela
tensão de cisalhamento em sua base; para inclinações pequenas $\sin\alpha \approx \alpha$ (o
teorema da laje com a inclinação da superfície).

\textbf{2.} $917\times9.81\times580\times0.0190 = \qty{99.1}{kPa}$;
$917\times9.81\times470\times0.0240 = \qty{101.5}{kPa}$; $917\times9.81
\times335\times0.0330 = \qty{99.4}{kPa}$.

\textbf{3.} A tensão é a mesma dentro de \qty{2}{\percent}, enquanto a
espessura quase se reduz à metade: a calota comporta-se quase como um material plástico, com $\tau_0 \approx
\qty{100}{kPa}$.

\textbf{4.} $2\times10^5/(917\times9.81) = \qty{22.23}{m}$.

\textbf{5.} $h = \sqrt{22.23\,(R - r)}$: \qty{577}{m}, \qty{472}{m} e
\qty{333}{m} a 15, 10 e \qty{5}{km} da margem, a poucos metros
dos valores medidos nas sondagens.

\textbf{6.} $\sqrt{22.23\times20000} = \qty{667}{m}$.

\textbf{7.} Como no \cref{exo:b3:glaciology:11}: $V = \int_0^R2\pi r
\sqrt{22.23(R - r)}\,\mathrm dr = \frac{8\pi}{15}\sqrt{2\tau_0/(\rho
g)}\,R^{5/2}$.

\textbf{8.} $\frac{8\pi}{15}\times4.715\times(\num{2e4})^{5/2} =
\qty{4.47e11}{m^3} = \qty{447}{km^3}$.

\textbf{9.} $\pi\times20^2 = \qty{1257}{km^2}$; espessura média
$447/1257 = \qty{0.356}{km}$, \qty{356}{m}, pouco mais da metade da
espessura central.

\textbf{10.} $2^{5/2} = 5.7$ vezes.

\textbf{11.} $A = \num{3.61e-13}\exp(-60000/(8.314\times263.15)) =
\qty{4.45e-25}{Pa^{-3}.s^{-1}}$, vezes \qty{3.156e7}{s}:
\qty{1.40e-17}{Pa^{-3}.yr^{-1}}.

\textbf{12.} A \qty{10}{km}, $H = \qty{471.5}{m}$, inclinação $471.5/20000 =
\num{0.0236}$, $\tau_b = \qty{100}{kPa}$; $u_s = \tfrac{2A}{4}\tau_b^3H =
0.5\times\num{1.40e-17}\times10^{15}\times471.5 = \qty{3.3}{m/yr}$.

\textbf{13.} $q = \tfrac{2A}{5}\tau_b^3H^2 = 0.4\times\num{1.40e-17}
\times10^{15}\times471.5^2 = \qty{1250}{m^2/yr}$; velocidade média $q/H =
\qty{2.6}{m/yr}$, quatro quintos da velocidade superficial.

\textbf{14.} $2\pi rq = \pi r^2\bar b$, logo $\bar b = 2q/r = 2\times1250/10^4
= \qty{0.25}{m}$ de gelo por ano sobre os \qty{10}{km} internos.

\textbf{15.} $h = 450$ onde $R - r = 450^2/22.23 = \qty{9.11}{km}$, em
$r = \qty{10.9}{km}$; a área de acumulação é o disco interno a ela, $(10.9/
20)^2 = 0.30$ da calota.

\textbf{16.} $A$ aumenta \num{10.2} vezes, e as velocidades sob as mesmas
tensões também. O gelo escoaria para fora mais depressa do que o balanço o
repõe: a calota afinaria e se espalharia até que as tensões, menores,
voltassem a igualar o balanço; uma calota plástica mais mole tem tensão de escoamento menor e
perfil mais achatado.

\textbf{17.} $\num{4.47e11}\times917 = \qty{4.10e14}{kg} = \qty{410}{Gt}$.

\textbf{18.} $410/360 = \qty{1.14}{mm}$.

\textbf{19.} $356/0.5 = 711$ anos.

\textbf{20.} À medida que a superfície desce, ela fica em ar mais quente e derrete mais depressa,
a retroalimentação altitude--balanço de massa; e, à medida que a calota afina, as tensões e o
escoamento também mudam. O afinamento acelera.

\textbf{21.} $410/243 = \num{1.7}\text{ anos}$: a calota inteira equivale a menos de dois
anos da perda da Groenlândia em 2010--2019.

\textbf{22.} A calota de Arkel tem um \omterm{def:b3:glaciology:sle}{equivalente em nível do mar} de
\textbf{\qty[text-series-to-math=true]{1.1}{mm}} (\qty{1.14}{mm}).
\end{solution}
